\section{Exact Optimization Methods}
\label{sec:exact-methods}

\subsection{Problem definition}
Let $G=(V,E)$ be a simple undirected graph with $n=|V|$ vertices and $m=|E|$ edges. A labeling is a bijection $f:V\rightarrow\{1,\ldots,n\}$. Its antibandwidth is the minimum label difference over all edges, and the Antibandwidth Problem (ABP) seeks a labeling that maximizes this minimum:
\begin{equation}
  \operatorname{AB}(G)
  = \max_{f\ \mathrm{bijective}}
    \min_{\{u,v\}\in E}|f(u)-f(v)|.
  \label{eq:abp-objective}
\end{equation}
This study evaluates four exact-solver backends for the same optimization objective: Gurobi MIP, CPLEX MIP, CPLEX Constraint Programming (CP), and OR-Tools CP-SAT. The formulations below describe the implementations in the project repository; they are not proposed as new ABP formulations.

\subsection{Common integer-label model}
For each vertex $v\in V$, let $x_v$ denote its assigned label, and let $K$ denote the minimum separation to be maximized. The common mathematical model is
\begin{align}
  \max\quad & K, \label{eq:exact-obj}\\
  \text{s.t.}\quad
    & 1\le x_v\le n && \forall v\in V, \label{eq:label-domain}\\
    & x_i\ne x_j && \forall\, i,j\in V,\ i<j, \label{eq:all-different}\\
    & |x_u-x_v|\ge K && \forall\,\{u,v\}\in E, \label{eq:edge-distance}\\
    & 1\le K\le n-1. \label{eq:k-domain}
\end{align}
The pairwise distinctness constraints in \eqref{eq:all-different}, together with the integer domains in \eqref{eq:label-domain}, make the labels a permutation of $\{1,\ldots,n\}$. The lower bound $K\ge 1$ is valid for an instance with at least one edge: every permutation assigns distinct labels to the endpoints of each edge. The CP-SAT implementation uses the slightly wider domain $0\le K\le n-1$; this does not change the optimum on the validated non-empty-edge instances, but the implementation detail is recorded here for reproducibility.

\subsection{Mixed-integer programming backends}
The Gurobi MIP and CPLEX MIP implementations use integer label variables $x_v$ and an integer objective variable $K$. They encode pairwise distinct labels and edge separation through binary orientation variables and indicator constraints.

For each unordered vertex pair $\{i,j\}$ with $i<j$, introduce $z_{ij}\in\{0,1\}$ and impose
\begin{align}
  z_{ij}=1 &\ \Longrightarrow\ x_i+1\le x_j,\\
  z_{ij}=0 &\ \Longrightarrow\ x_j+1\le x_i.
\end{align}
Exactly one of the two binary cases is active, so every pair of vertices receives distinct labels. Since all labels are integers in $\{1,\ldots,n\}$, the resulting assignment is a permutation.

For every edge $\{u,v\}\in E$, introduce a binary variable $y_{uv}\in\{0,1\}$ and impose
\begin{align}
  y_{uv}=1 &\ \Longrightarrow\ x_u-x_v\ge K,\\
  y_{uv}=0 &\ \Longrightarrow\ x_v-x_u\ge K.
\end{align}
Thus, the selected orientation enforces the required absolute separation on each edge. Both MIP backends implement these implications as native indicator constraints; they do not use manually selected big-$M$ constants.

The resulting structural variable count is
\[
  n+1+\binom{n}{2}+m,
\]
comprising $n$ label variables, one objective variable, one orientation variable per vertex pair, and one direction variable per edge. The implementation contains two indicator implications per vertex pair and two per edge. These are structural counts of the submitted model, not a claim about the internal presolved model generated by either solver.

\subsection{Constraint-programming backends}
The CPLEX CP implementation represents the labels as integer variables $x_v\in\{1,\ldots,n\}$, posts a global \texttt{all\_diff} constraint over all labels, and adds
\[
  |x_u-x_v|\ge K
  \qquad \forall\,\{u,v\}\in E.
\]
It maximizes $K$. At the modeling level this uses $n+1$ integer variables (the $n$ labels and $K$), one global all-different constraint, and one absolute-distance constraint per edge.

The CP-SAT implementation also uses integer label variables and an all-different constraint. For each edge $\{u,v\}$, it introduces an integer distance variable $d_{uv}$ and posts
\begin{align}
  d_{uv} &= |x_u-x_v|,\\
  K &\le d_{uv}.
\end{align}
It maximizes $K$. Its structural variable count is $n+1+m$. CP-SAT is configured with one search worker. As with the CP formulation, global and absolute-value constraints may be internally transformed by the solver; the structural count should not be interpreted as a direct measure of post-presolve memory.

\subsection{Degree-based upper bound}
Let $\Delta=\max_{v\in V}\deg(v)$. For every graph with at least one edge, a valid upper bound is
\begin{equation}
  \operatorname{AB}(G)\le n-\Delta.
  \label{eq:degree-upper-bound}
\end{equation}
To see this, choose a vertex $w$ of degree $\Delta$ and fix its label $x_w$. If a labeling has antibandwidth at least $K$, each of the $\Delta$ neighbours of $w$ must receive a label whose distance from $x_w$ is at least $K$. Among the $n$ available labels, at most $n-K$ labels can be at such a distance from a fixed label. Hence $\Delta\le n-K$, which gives $K\le n-\Delta$.

The bound is recorded in the benchmark manifest and used as an external consistency bound. It is not, by itself, an optimality certificate. In the current Exact implementations, the model domain for $K$ remains the natural range up to $n-1$; the degree bound is not passed into each Exact model. In particular, a run is classified as optimal only when the solver reports a proof of optimality and the incumbent passes the independent verifier.

\subsection{Solution verification and termination status}
A solver status is not accepted as a certificate of a valid labeling. Every incumbent is decoded into a mapping from all vertices $0,\ldots,n-1$ to labels $1,\ldots,n$ and is independently checked by the shared verifier. The verifier requires a complete bijection and recomputes the objective as
\[
  \min_{\{u,v\}\in E}|x_u-x_v|.
\]
The recomputed value must agree with the solver-reported objective before the result is retained.

The project uses the following common result semantics:
\begin{itemize}
  \item \texttt{OPT}: a verified incumbent and a solver proof of optimality; the reported lower and upper bounds coincide with the verified objective.
  \item \texttt{FEAS}: a verified incumbent exists, but optimality has not been proved within the run budget.
  \item \texttt{TO}: the time budget expires without a verified incumbent.
  \item \texttt{MO}: a memory-limit event is established by the solver or execution supervisor. Any valid incumbent and its bound information must be preserved when available.
  \item \texttt{ERR}: an execution, license, backend, malformed-output, invalid-labeling, or objective-consistency error prevents the run from being accepted. A license failure is not reclassified as a timeout.
\end{itemize}
For non-optimal feasible runs, the verified objective is retained as the lower bound. An upper bound is taken from the solver when available and otherwise from a valid safe bound. An unverified labeling is never promoted to an incumbent.

\paragraph{Implementation-conformance note.}
The status definitions above describe the intended project-level contract. A source audit found that the current adapters do not yet normalize every termination cause identically: for example, the tested CPLEX CP path returns \texttt{FEAS} when a verified incumbent exists at a memory stop, whereas the MIP adapters can return \texttt{MO} for a recognized memory-limit termination even when an incumbent is available. The CPLEX CP extraction path also requires a distinction between the absence of an incumbent and an incumbent that fails decoding or verification. These cases are tracked as open conformance items in the P1 Exact methodology review and must be resolved, or explicitly specified and handled in the analysis, before the status comparison is frozen.

\subsection{Time accounting and implementation scope}
Each Exact backend measures elapsed time from the beginning of its solver call, including backend import and model construction. The optimizer receives only the remaining portion of the configured whole-run budget after model construction. If the budget has already expired during construction, optimization is not started. The runner additionally records process-level \texttt{wall\_time}, which includes worker startup and input loading and is the primary end-to-end runtime used in the experimental comparison. Solver-reported \texttt{t\_encode}, \texttt{t\_solve}, and \texttt{t\_total} are retained as diagnostic measurements and are not mixed with \texttt{wall\_time} in one runtime or PAR2 comparison.

A brute-force oracle enumerates permutations and evaluates each labeling using the same verifier. It is restricted to small instances (default maximum $n=9$) and is used for correctness tests, not for the Harwell--Boeing benchmark runs.

\subsection{Relation to prior exact approaches}
Exact optimization of the ABP has been studied using mixed-integer programming and constraint programming. Sinnl reports MIP formulations, branch-and-cut and iterative solution approaches, as well as a CP formulation, and evaluates them on the Harwell--Boeing instances \citep{Sinnl2021}. The present implementation uses direct integer label variables with indicator-based pairwise ordering in its MIP backends, and global all-different/absolute-distance constraints in its CP backends. Its purpose in this project is to provide consistent Exact baselines for a controlled comparison with the selected SAT configuration, rather than to claim a new formulation or to reproduce the stronger algorithms and valid inequalities proposed in prior work.
